\documentclass[12pt,a4paper]{article}
\usepackage[french]{babel}
\usepackage[T1]{fontenc}
\usepackage[utf8]{inputenc}
\usepackage{amsmath, amssymb}
\usepackage{geometry}
\geometry{top=1cm, bottom=2cm, left=1cm, right=1cm}
\usepackage{multicol}
\usepackage{graphicx}
\usepackage{tcolorbox}
\usepackage{multicol}

\begin{document}

\centering\textbf{Devoir Surveillé de Mathématiques - Proportions et Évolutions}\\
\vspace{0.5cm}

\noindent\textbf{Nom :} \underline{\hspace{5cm}} \hfill \textbf{Prénom :} \underline{\hspace{5cm}} \\
\textbf{Classe :} \underline{\hspace{5cm}}\\
\vspace{1cm}

\centering\textit{\textbf{Chaque question vaut 1 point, sauf les questions avec calculs complexes pouvant être notées jusqu’à 3 points selon la qualité de la rédaction. La présentation est évaluée sur 2 points. Calculatrice autorisée.}}

\section*{Partie 1 : Proportions}

\subsection*{Exercice 1 : Étude d’une population}
Dans un lycée de 500 élèves, une classe représente 7\% des élèves du lycée. Dans cette classe :
\begin{itemize}
    \item 14 élèves étudient uniquement l’anglais,
    \item 12 élèves étudient uniquement l’espagnol,
    \item 5 élèves étudient les deux langues.
\end{itemize}

\begin{enumerate}
    \item Quelle est la proportion d'élèves qui étudient l’espagnol par rapport à la classe ?\\
    \begin{tcolorbox}[colframe=black, colback=white, width=\linewidth, height=3cm]
    \end{tcolorbox}

    \item Quelle est la proportion d’élèves qui étudient l’anglais parmi ceux qui étudient aussi l’espagnol ?\\
    \begin{tcolorbox}[colframe=black, colback=white, width=\linewidth, height=3cm]
    \end{tcolorbox}

    \item Quelle est la proportion d’élèves qui étudient l’anglais et l’espagnol par rapport à la classe ? Donner deux méthodes :
    \begin{itemize}
        \item À partir de l’énoncé
        \item À partir des questions précédentes
    \end{itemize}
    \begin{tcolorbox}[colframe=black, colback=white, width=\linewidth, height=5cm]
    \end{tcolorbox}

    \item Quelle est la proportion d’élèves qui n’étudient aucune des deux langues ?\\
    \begin{tcolorbox}[colframe=black, colback=white, width=\linewidth, height=3cm]
    \end{tcolorbox}
\end{enumerate}

\section*{Partie 2 : Évolutions et coefficients multiplicateurs}

\subsection*{Exercice 2 : Réduction et augmentation}

Du plâtre pour moulures est vendu dans un magasin. Kévin applique une réduction de 20\%, puis deux semaines plus tard une augmentation de 20\%.

\begin{enumerate}
    \item Cette augmentation permet-elle de retrouver le prix initial ? Justifier par un calcul.\\
    \begin{tcolorbox}[colframe=black, colback=white, width=\linewidth, height=4cm]
    \end{tcolorbox}

    \item Inès, sa collègue, affirme qu’il faut multiplier le prix réduit par 1{,}25 pour retrouver le prix initial. A-t-elle raison ? Justifier par un calcul.\\
    \begin{tcolorbox}[colframe=black, colback=white, width=\linewidth, height=4cm]
    \end{tcolorbox}
\end{enumerate}

\subsection*{Exercice 3 : Coefficient multiplicateur global}

\begin{enumerate}
    \item Déterminer le coefficient multiplicateur global dans les cas suivants :\\
    \begin{itemize}
        \item une réduction de 20\% suivie d’une augmentation de 10\% ;
        \item une réduction de 25\% suivie d’une augmentation de 50\%.
    \end{itemize}
    \begin{tcolorbox}[colframe=black, colback=white, width=\linewidth, height=5cm]
    \end{tcolorbox}

    \item En déduire, dans chaque cas, le taux d’évolution global.\\
    \textit{Ainsi cela correspond à une \dotfill de \dotfill \%.}\\
    \begin{tcolorbox}[colframe=black, colback=white, width=\linewidth, height=3cm]
    \end{tcolorbox}
\end{enumerate}

\subsection*{Exercice 4 : Retrouver un prix initial}

\begin{enumerate}
    \item Un téléphone a subi une réduction de 30\%. Son prix après réduction est de 811{,}30 €. Quel était son prix initial ?\\
    \begin{tcolorbox}[colframe=black, colback=white, width=\linewidth, height=4cm]
    \end{tcolorbox}

    \item Un ordinateur a subi une réduction de 30\% suivie d’une augmentation de 20\%. Son prix final est de 504 €. Quel était son prix de départ ?\\
    \begin{tcolorbox}[colframe=black, colback=white, width=\linewidth, height=4cm]
    \end{tcolorbox}
\end{enumerate}

\end{document}
